% Generated draft: requires mathematical and source review.
\documentclass[11pt]{article}
\usepackage[margin=25mm]{geometry}
\usepackage{fontspec}
\setmainfont{Latin Modern Roman}
\usepackage{amsmath,amssymb,mathtools,unicode-math}
\setmathfont{Latin Modern Math}
\usepackage{enumitem,needspace,xurl,hyperref}
\setlength{\parindent}{0pt}
\setlength{\parskip}{5pt}
\setlength{\emergencystretch}{3em}
\newcommand{\N}{\mathbb{N}}
\newcommand{\Z}{\mathbb{Z}}
\newcommand{\Q}{\mathbb{Q}}
\newcommand{\R}{\mathbb{R}}
\newcommand{\C}{\mathbb{C}}
\newcommand{\F}{\mathbb{F}}
\newcommand{\A}{\mathbb{A}}
\newcommand{\PP}{\mathbb{P}}
\newcommand{\E}{\mathbb{E}}
\newcommand{\eps}{\varepsilon}
\newcommand{\dd}{\,\mathrm{d}}
\begin{document}
\section*{UnsolvedMath ID: 20001638}
\textbf{Source problem code:} catalogue entry
\subsection{Definitions and mathematical statement}
This canonical record is editorial context, not Problem 1.1. The supplied verification says the original paragraph described the \emph{development} of Problem 1.1 as intermingled with discussion, then introduced candidate complexes. The actual five-part problem is assigned to \texttt{AIM-GEOMETRIC\_GROUP\_THEORY-0148}, and the candidate audit to \texttt{AIM-GEOMETRIC\_GROUP\_THEORY-0149}. Their full questions are not present here and are not silently substituted for this record.

The aid separately develops a comparison criterion. For a group \(\Gamma\) acting isometrically on a nonempty connected metric graph \(X\), define stable translation length
\[
 \tau_X(g)=\lim_{m\to\infty}\frac{d_X(x,g^m x)}m,
\]
independent of \(x\). Let \(X,Y\) be two such \(\Gamma\)-graphs, and suppose \(f:X\to Y\) satisfies
\[
 d_Y(fx,fx')\leq L_fd_X(x,x')+C_f,\qquad
 d_Y(f(gx),g f(x))\leq E_f
\]
for all \(x,x'\) and \(g\in\Gamma\), with constants independent of all those arguments. Then the reported directional conclusion is
\[
 \tau_Y(g)\leq L_f\tau_X(g)\quad(g\in\Gamma).
\]
If also \(h:Y\to X\) has the same type of uniform estimates with constant \(L_h\), the two inequalities compare translation lengths in both directions. The maps need not be quasi-inverses. Uniform coarse equivariance for all group elements is material; a generatorwise error accumulating with word length would not give the same limit inequality.

For the aid's application, take \(\Gamma=\operatorname{Out}(F_n)\), \(n\geq3\), and \(Y=\mathcal F_n^{(1)}\), the one-skeleton of the free-factor order complex. Its vertices are conjugacy classes of proper nontrivial free factors and its simplices are chains under conjugate inclusion. The reported benchmark is positive stable length for fully irreducible outer automorphisms and zero stable length for polynomially growing ones. Fully irreducible means no positive power preserves a proper nontrivial free-factor conjugacy class. Thus a uniform coarse-Lipschitz coarsely equivariant map
\(X\to\mathcal F_n^{(1)}\)
transfers the positive-length obstruction to \(X\), whereas a map
\(\mathcal F_n^{(1)}\to X\)
transfers zero length. Both maps give both properties. They do not on their own prove hyperbolicity, boundary descriptions or link-connectivity conditions for \(X\).

The independent link audit uses the following exact recursion. At a vertex \([A]\) of the free-factor order complex,
\[
 \operatorname{Lk}([A])=
 \Delta(\mathcal P_{<[A]})*\Delta(\mathcal P_{>[A]}),\qquad
 \Delta(\mathcal P_{<[A]})\cong\mathcal F(A),
\]
where \(\mathcal P\) is the free-factor poset, \(\Delta\) means its order complex, and \(*\) is simplicial join. The descending factor consists exactly of conjugacy classes of proper nontrivial free factors of \(A\); when \(A\) has rank one it is empty. No identification of the ascending factor with an unrelated absolute free-factor complex is asserted. Consequently this lower-link recursion and the needed connectivity of relevant links must be checked separately from translation-length comparisons. These are reported mathematical supplements to an editorial fragment, not the missing five-part original statement.
\subsection{Short English statement}
The record is an editorial bridge. The aid adds a precise two-direction comparison for stable translation lengths and a separate free-factor lower-link recursion, while preserving the missing original problem and candidate audit as different records.
\subsection{Sources}
\begin{itemize}
\item Local registry: permanent UnsolvedMath ID \(20001638\), exact public code \texttt{catalogue entry}; \texttt{aim-geometric-group-theory-notes.json}, \texttt{notes[146]}.
\item AIM workshop ``The geometry of the outer automorphism group of a free group'' (2010), Section 1 ``Curve complex analogues'': \url{https://aimath.org/WWN/outerauto/outerauto.pdf}. The supplied verification identifies this as a damaged introductory paragraph and locates actual Problem 1.1 in public record 0148, followed by candidate audit 0149. The registry marks the statement \texttt{unrecoverable}. No proposer or literature entry is supplied.
\item Research attempt 1: \texttt{context\_only}, status at run \texttt{not\_a\_problem}, rigor label \texttt{proved}, novelty confidence low, suggested difficulty L2, classification \texttt{OPEN-TRIAGE}. It reports the coarse-comparison criterion and lower-link recursion as a candidate reduction. No authoritative separate references for benchmark dynamics or link facts are supplied. The current registry status is open, as metadata for the fragment.
\end{itemize}
\end{document}
